\section{Results and external comparisons}
\label{sec:results}

The central inference is the simultaneous shared-$\eps^2$ limit from the product
likelihood in Eq.~\eqref{eq:combined-likelihood}. Individual data-set limits and
local $p_0$ scans are important diagnostics, but they are not the combination
prescription. At fixed mass, the common coupling predicts different expected yields in
the active data samples through their own $K_d(m)$ factors, and the likelihood profiles
the data-set-specific GPR covariance nuisance parameters independently.

Figure~\ref{fig:combined-95} shows the corrected staged 95\% confidence-level result
from the 2015, 2016 10\%, and 2021 1\% inputs. The observed curve and the expected
bands are shown on the same $\eps^2$ axis. Because this result uses partial 2016 and
2021 samples, it should be read as the current internal-review baseline rather than
the final full-statistics exclusion.

\begin{figure*}[t]
\centering
\ifigure{0.82\textwidth}{combined_limit_95cl.png}
\caption{Corrected simultaneous 2015+2016 10\%+2021 1\% upper limit on $\eps^2$ at
95\% confidence level. The green and yellow bands are the corrected background-only
expected intervals, and the black curve is the observed shared-coupling limit. This is
the current staged GPR result, not a final full-statistics exclusion.}
\label{fig:combined-95}
\end{figure*}

For external visible-dark-photon comparisons, a 90\% $\CLs$ presentation is also
prepared. Figure~\ref{fig:combined-90} shows the corrected 90\% result with the
minimal-model dimuon-threshold branching correction included in the coupling
interpretation. Below $2m_\mu$, this correction is unity; above threshold it weakens
the electron-channel coupling interpretation because additional visible decay channels
contribute to the total width.

\begin{figure*}[t]
\centering
\ifigure{0.82\textwidth}{combined_90cl_dimuon.png}
\caption{Corrected simultaneous 90\% $\CLs$ upper limit on $\eps^2$ with the
minimal-model dimuon-threshold branching correction. This figure is included to support
like-for-like external-contour comparisons. The final PRD submission should freeze the
confidence-level convention and systematic prescription before quoting the exclusion.}
\label{fig:combined-90}
\end{figure*}

Observed combined limits are not expected to be everywhere stronger than the best
individual observed curve. An upward fluctuation in one active sample can weaken a
shared-coupling limit, while a downward fluctuation in one sample can make an
individual observed curve appear stronger. Figure~\ref{fig:individual-overlay} is
therefore shown as a diagnostic of the combined fit, not as a way to define the
combined result.

\begin{figure*}[t]
\centering
\ifigure{0.82\textwidth}{individual_combined_overlay.png}
\caption{Observed individual and simultaneous 95\% confidence-level upper limits on
$\eps^2$. The simultaneous curve is obtained from the joint profiled likelihood. The
individual curves identify which data samples are active and which observed
fluctuations influence the shared-coupling result.}
\label{fig:individual-overlay}
\end{figure*}

The discovery-oriented diagnostics are shown in Fig.~\ref{fig:combined-discovery}.
The local $p_0$ curve is the fixed-mass one-sided discovery test. The corresponding
global overlay is an effective-trials approximation included for review. It should be
replaced by the scan-level toy calibration of $q_{0,\max}$ before any final discovery
claim is made.

\begin{figure*}[t]
\centering
\begin{minipage}{0.48\textwidth}
\centering
\ifigure{\linewidth}{combined_p0.png}\\
(a) Local $p_0$ with approximate global guide.
\end{minipage}
\hfill
\begin{minipage}{0.48\textwidth}
\centering
\ifigure{\linewidth}{combined_Z.png}\\
(b) Gaussian-equivalent local significance.
\end{minipage}
\caption{Discovery diagnostics for the staged simultaneous analysis. These curves test
for upward local excesses and are distinct from the observed-limit tail diagnostics.
The global references are effective-trials guides pending the final scan-level toy
calibration.}
\label{fig:combined-discovery}
\end{figure*}

The current full-statistics curves are projections. They scale the partial 2016 and
2021 inputs to their anticipated full exposures while holding the 2015 sample fixed.
Such curves are useful for judging the expected reach of the GPR strategy and for
designing the final paper figures, but they are not observed full-statistics
exclusions. Figure~\ref{fig:projected-phase-space} shows the projected full-data GPR
reach in the same visible-dark-photon phase space as Fig.~\ref{fig:intro-context}.

\begin{figure*}[t]
\centering
\ifigure{0.82\textwidth}{projected_phase_space_context.png}
\caption{Projected full-data HPS GPR reach in visible-dark-photon parameter space. The
world-data contours and HPS engineering-run contours are shown for orientation. The
thick projected GPR curve is an observed-equivalent projection from the current staged
inputs and should not be described as an observed exclusion until the full data samples
and final systematic model are frozen.}
\label{fig:projected-phase-space}
\end{figure*}

\clearpage
